\section{Results}
\label{sec:results}

\subsection{Primary outcome}

No primary chain completed before the results freeze. This paper reports the
preregistered design, the validated scaffold, the positive-control status, and
the validity tooling. C-002 and C-003 remain untested. We make no empirical
claim about allocation policy.

The three primary configurations remain unfrozen, and the chain runner declines to launch
an unfrozen configuration. This is a design decision awaiting resolution, not a compute
shortfall: no accelerator time would advance it. Reporting the absence directly is the
precommitted response to this outcome.

\begin{table}[h]
\centering
\caption{Primary chain-level results. This file is generated by the analysis pipeline from
immutable chain artifacts and is empty while no primary chain exists.}
% AUTO-GENERATED by scripts/generate_pilot_outputs.py. DO NOT EDIT VALUES BY HAND.
% Derived from chain_result.json artifacts of primary_pilot_v2_2026-08-20.
% Both budget axes hold, so the preregistered contrast is reported below the
% table. Strongest eligible non-joint baseline: selection_only.
\begin{tabular}{lrrrrr}
\toprule
Policy & Chains & Human tokens & AUC regret & SD & Tail retention \\
\midrule
No rescue (control) & 5 & 0 & 2.6431 & 0.0140 & 0.8676 \\
Fresh random & 5 & 749,869 & 2.5364 & 0.0080 & 0.8769 \\
Schedule-only & 5 & 749,878 & 2.5261 & 0.0194 & 0.8846 \\
Selection-only & 5 & 749,827 & 2.2931 & 0.0250 & 0.9024 \\
Joint time-and-mode & 5 & 749,827 & 2.3034 & 0.0221 & 0.9024 \\
\bottomrule
\end{tabular}

% Primary contrast, joint minus selection_only (paired by seed, 95\% t interval):
%   +0.01026 [-0.00916, +0.02968], +0.45\% relative.
%   Practical-equivalence region is +/-0.05073 (2\% of the fresh-random mean, the denominator U-006 was settled against).
%   Interval lies entirely inside that region: yes.
% Budget axes: human spread 0.0381\%, total spread 0.0000\%, permitted 0.2000\%.

\end{table}

\subsection{Positive control}

The published positive control reproduces qualitatively. Both arms complete every planned
generation; the expected arm ordering and degradation direction are recovered; a clean
rerun on independent hardware reaches the same conclusion, with every reported quantity
agreeing to well under one per cent across the two executions; all deviations are
enumerated; and
the reported table is regenerated from saved metrics by the same command that produced
them, with every recorded artifact hash verified.

Agreement with the published endpoint is assessed against a five per cent relative difference
adopted as an internal engineering tolerance. It is not proof of exact replication, and no
downstream claim rests on it.

This establishes that the pipeline reproduces a known result on real assets. It says
nothing about allocation policy, which is the question this paper poses and does not
answer.
